% ============================================================
% SECTION 1: OPENING & CONTEXT (Slides 1–4)
% ============================================================

\section{Context}

% ----------------------------------------------------------
% Slide 2: Executive Summary
% ----------------------------------------------------------
\begin{frame}{Executive Summary}
  \begin{itemize}
    \item Evaluated VSP-LLM lip-reading system on \textbf{1,497 real-world YouTube segments}
    \item Standard metric (WER): \werbad{64.1\%} --- \textbf{2.5$\times$} worse than benchmark
    \item Our new \keyword{Intelligibility Score} reveals \textbf{40\%} of output is actually useful ---
          \keyword{3.5$\times$} more than WER admits
    \item Built production system: 8-stage pipeline, standalone container, 37 bugs fixed
    \item Produced \textbf{8 research reports} covering evaluation, tuning, and methodology
    \item Clear improvement roadmap targeting \keyword{27--42\% WER}
  \end{itemize}

  \vspace{6pt}
  \centering
  {\textcolor{ArgosTeal}{\textit{``Half the output delivers useful information.\\
  The starting point is much better than WER says.''}}}

  \note{
    BOTTOM LINE UP FRONT.
    Three months of research and engineering. Key finding: WER dramatically overstates
    failure. Our Intelligibility Score shows 40\% of output is properly captured, not the
    11\% WER suggests. Complete production system delivered. Clear roadmap to improve.
  }
\end{frame}

% ----------------------------------------------------------
% Slide 3: What Is Lip Reading?
% ----------------------------------------------------------
\begin{frame}{What Is Lip Reading?}
  \twocol{0.55}{
    \begin{itemize}
      \item Speech recognition from \keyword{visual input only}
      \item No audio signal --- mouth movements $\rightarrow$ text
      \item 50--70\% of English sounds are \alert{invisible} on lips
            \notetext{(p/b/m identical, t/d/n/s/z/l identical)}
      \item Applications: surveillance, accessibility, noisy environments
    \end{itemize}
    \vspace{6pt}
    \notetext{Video: 33-word insurance policy perfectly transcribed (WER 0\%)}
  }{0.42}{
    \centering
    \videocaptioned{\linewidth}%
      {IEa7qEkMvfQ_3__c5447488_thumb.png}%
      {videos/source/IEa7qEkMvfQ_3__c5447488_with_hyp.mp4}%
      {Perfect transcription --- WER 0\%, IS 5.0}
  }

  \note{
    Play the video FIRST, explain after.
    33 words about insurance policies, perfectly lip-read.
    Shows the technology fundamentally works on complex domain-specific content.
    ``This is what it looks like when it works. Now let's see what happens on real-world data.''
  }
\end{frame}

% ----------------------------------------------------------
% The Fundamental Challenge: Visemes
% ----------------------------------------------------------
\begin{frame}{The Fundamental Challenge: Visemes}
  \twocol{0.52}{
    {\small 50--70\% of English sounds are \alert{invisible} on lips.\\
    Multiple sounds produce identical mouth shapes:}

    \vspace{6pt}

    \begin{tabular}{@{}ll@{}}
      \toprule
      \textbf{Viseme Group} & \textbf{Sounds} \\
      \midrule
      Bilabial & p, b, m \\
      Alveolar & t, d, n, s, z, l \\
      Velar & k, g, ng \\
      Labiodental & f, v \\
      \bottomrule
    \end{tabular}

    \vspace{6pt}
    \notetext{``Viseme'' = visual phoneme. One viseme maps\\to multiple phonemes.}
  }{0.45}{
    {\small\textbf{Confusable word pairs:}}

    \vspace{4pt}

    \begin{tabular}{@{}l@{$\;\leftrightarrow\;$}l@{}}
      \toprule
      \multicolumn{2}{c}{\textbf{Identical on Lips}} \\
      \midrule
      admiral & animal \\
      mom & bomb \\
      pat & bat \\
      collar & color \\
      \bottomrule
    \end{tabular}

    \vspace{8pt}
    \centering
    {\textcolor{ArgosTeal}{\textbf{Context is the ONLY\\disambiguation signal ---\\this is why the LLM matters.}}}
  }

  \note{
    50--70\% of sounds are invisible on lips --- same mouth shape for multiple phonemes.
    ``admiral'' vs ``animal'': identical bilabial + alveolar sequence.
    This is the fundamental challenge that makes lip-reading so hard.
    Context (the LLM component) is the ONLY way to disambiguate.
    This motivates the entire architecture: a frozen visual encoder + a powerful language model.
  }
\end{frame}

% ----------------------------------------------------------
% Slide 3: Model Architecture
% ----------------------------------------------------------
\begin{frame}[shrink=5]{Model Architecture}
  \begin{center}
    % Model architecture diagram: Video Frames → AV-HuBERT → Projection → LLaMA-2-7B → Text
% Redesigned for clarity: distinct shapes, white text on color, visible badges
\begin{tikzpicture}[
  % ── Data nodes (input/output) ────────────────────────────────────────
  data/.style={
    draw=ArgosGray, fill=ArgosGray!15, line width=1pt,
    rounded corners=3pt,
    minimum height=1.1cm, minimum width=1.7cm,
    align=center, inner sep=5pt, text=ArgosNavy,
    font=\small\bfseries
  },
  % ── Processing component boxes ────────────────────────────────────────
  frozen/.style={
    draw=ArgosTeal, fill=ArgosTeal, line width=1.2pt,
    rounded corners=6pt,
    minimum height=1.4cm, minimum width=2.6cm,
    align=center, inner sep=6pt, text=white,
    font=\small\bfseries
  },
  bridge/.style={
    draw=ArgosGray!60, fill=ArgosGray!20, line width=1pt,
    rounded corners=4pt,
    minimum height=1.0cm, minimum width=1.4cm,
    align=center, inner sep=4pt, text=ArgosNavy,
    font=\scriptsize\bfseries
  },
  llm/.style={
    draw=ArgosDarkBlue, fill=ArgosDarkBlue, line width=1.2pt,
    rounded corners=6pt,
    minimum height=1.4cm, minimum width=2.6cm,
    align=center, inner sep=6pt, text=white,
    font=\small\bfseries
  },
  % ── Status pill badge ────────────────────────────────────────────────
  pill/.style={
    rounded corners=3pt, inner sep=2.5pt,
    font=\tiny\bfseries, text=white, fill=#1
  },
  % ── Dimension annotation ─────────────────────────────────────────────
  dim/.style={font=\tiny, ArgosGray, align=center},
  % ── Arrow ────────────────────────────────────────────────────────────
  arr/.style={-{Stealth[length=3mm,width=2mm]}, line width=1pt, ArgosGray},
  node distance=0.55cm
]

% ── Nodes ────────────────────────────────────────────────────────────

\node[data] (input) {Video\\Frames};

\node[frozen, right=0.6cm of input] (hubert)
  {AV-HuBERT\\[2pt]{\scriptsize Visual Encoder}};

\node[bridge, right=0.5cm of hubert] (proj)
  {Linear\\Proj.};

\node[llm, right=0.5cm of proj] (llama)
  {LLaMA-2-7B\\[2pt]{\scriptsize Language Model}};

\node[data, right=0.6cm of llama] (output) {Text\\Output};

% ── Arrows ───────────────────────────────────────────────────────────

\draw[arr] (input.east)  -- (hubert.west);
\draw[arr] (hubert.east) -- (proj.west);
\draw[arr] (proj.east)   -- (llama.west);
\draw[arr] (llama.east)  -- (output.west);

% ── Dimension labels (above arrows) ──────────────────────────────────

\node[dim] at ($(hubert.east)!0.5!(proj.west)+(0,0.52)$)
  {1024-dim};
\node[dim] at ($(proj.east)!0.5!(llama.west)+(0,0.52)$)
  {1024{\scriptsize$\to$}4096};
\node[dim] at ($(llama.east)!0.5!(output.west)+(0,0.52)$)
  {4096-dim};

% ── Status badges (below each component) ─────────────────────────────

\node[pill=ArgosGray!70, below=3pt of hubert]
  {\faIcon{lock}~Frozen};

\node[pill=ArgosOrange, below=3pt of proj]
  {Trained};

\node[pill=ArgosGray!70, below=3pt of llama]
  {\faIcon{lock}~Frozen};
\node[pill=ArgosDarkBlue!80, below=14pt of llama]
  {QLoRA r=16};

\end{tikzpicture}
  \end{center}

  \vspace{2pt}
  \begin{columns}[T]
    \begin{column}{0.32\textwidth}
      \centering
      {\small\keyword{Visual Encoder}}\\
      {\scriptsize AV-HuBERT (frozen)\\1024-dim features}
    \end{column}
    \begin{column}{0.32\textwidth}
      \centering
      {\small\keyword{Projection}}\\
      {\scriptsize Linear 1024$\rightarrow$4096\\Trainable bridge}
    \end{column}
    \begin{column}{0.32\textwidth}
      \centering
      {\small\keyword{Language Model}}\\
      {\scriptsize LLaMA-2-7B (4-bit QLoRA)\\r=16, 12.6M params (0.19\%)}
    \end{column}
  \end{columns}

  \vspace{4pt}
  \notetext{LLM is swappable --- Llama 3.1 8B has same hidden\_size (4096). Only LoRA adapters + projection are trained.}

  \note{
    Three blocks: frozen visual encoder, trainable projection, quantized LLM.
    Key point: visual encoder is FROZEN --- all learning happens in projection + LoRA.
    LLM swap is trivial: same hidden dimension, 1--2 hours setup.
    Only 0.19\% of parameters are trained --- very parameter-efficient.
  }
\end{frame}

% ----------------------------------------------------------
% How It Works: From Video to Text
% ----------------------------------------------------------
\begin{frame}{How It Works: From Video to Text}
  \vspace{2pt}
  {\small Five steps from raw video to transcribed text:}

  \vspace{6pt}

  \begin{enumerate}
    \item \keyword{Video Frames} --- raw video at 25 fps
    \item \keyword{Mouth Crop} --- 96$\times$96 grayscale, face detection + alignment
    \item \keyword{Visual Features} --- AV-HuBERT encoder, 1024-dim vectors per frame
    \item \keyword{Projection} --- linear layer maps 1024 $\rightarrow$ 4096 dimensions
    \item \keyword{LLM Generates Text} --- Llama-2-7B decodes visual tokens into English
  \end{enumerate}

  \vspace{6pt}

  \begin{center}
  \begin{tikzpicture}[
    box/.style={draw=ArgosTeal!60, fill=ArgosTeal!8, thick,
                rounded corners=2pt, minimum height=0.7cm,
                font=\scriptsize, align=center},
    arr/.style={-{Stealth[length=2mm]}, ArgosTeal, thick},
    node distance=0.15cm
  ]
    \node[box] (v) {Video\\25 fps};
    \node[box, right=of v] (c) {Mouth Crop\\96$\times$96};
    \node[box, right=of c] (f) {AV-HuBERT\\1024-dim};
    \node[box, right=of f] (p) {Projection\\4096-dim};
    \node[box, right=of p] (l) {LLM\\Text Output};

    \draw[arr] (v) -- (c);
    \draw[arr] (c) -- (f);
    \draw[arr] (f) -- (p);
    \draw[arr] (p) -- (l);
  \end{tikzpicture}
  \end{center}

  \vspace{4pt}
  \centering
  \notetext{Visual encoder is \textbf{FROZEN} --- all learning in projection + LoRA adapters (0.19\% of params).}

  \note{
    Concrete walkthrough for non-specialists. Bridges the abstract architecture to a step-by-step process.
    Key: the visual encoder (step 3) is frozen --- it was pre-trained on 30K hours of video.
    All learning happens in the projection layer and LoRA adapters: only 12.6M of 6.7B parameters.
    This is why the visual encoder is the bottleneck --- we cannot improve its representations without
    unfreezing and retraining it, which requires massive data.
  }
\end{frame}

% ----------------------------------------------------------
% Slide 4: The Benchmark
% ----------------------------------------------------------
\begin{frame}[shrink=5]{The Benchmark}
  \twocol{0.50}{
    \begin{block}{Published Result}
      \keyword{25.4\% WER} on LRS3\\[2pt]
      {\small 1,000+ hours of TED talks\\
      Controlled studio lighting\\
      Frontal face, professional speakers\\
      Clean audio transcription labels}
    \end{block}

    \vspace{4pt}

    \begin{alertblock}{Our Dataset}
      {\small \textbf{1,497 segments} from 100+ diverse\\
      YouTube videos --- cooking shows, news,\\
      tutorials, vlogs, interviews\\[2pt]
      No curation, no controlled lighting,\\
      no frontal face requirement}
    \end{alertblock}

    \vspace{4pt}

    \begin{alertblock}{Our Questions}
      \begin{enumerate}
        \item How does this perform on \alert{real-world video}?
        \item Is \alert{WER even the right metric}?
      \end{enumerate}
    \end{alertblock}
  }{0.47}{
    \halfwidthfig{P2_paper_vs_reality.png}
  }

  \note{
    Paper reports 25.4\% WER on LRS3 benchmark --- ``easy mode.''
    LRS3: 1,000+ hours TED talks, studio-quality, frontal face, professional speakers.
    Our dataset: 1,497 segments from 100+ diverse YouTube videos.
    No curation, no controlled lighting, no frontal face requirement.
    Two questions frame the entire presentation:
    1. Real-world performance (spoiler: 2.5x worse)
    2. WER as a metric (spoiler: dramatically overstates failure)
  }
\end{frame}
